\chapter{Resultados}

% Apresente primeiro os resultados e depois sua interpretação.
Organize os achados de acordo com os objetivos e as hipóteses da pesquisa.

\section{Análise descritiva}

Caracterize a amostra por meio de estatísticas, tabelas e gráficos. Numere e
identifique todas as figuras e tabelas, cite suas fontes e faça referência a
elas no texto.

\section{Resultados principais}

Apresente as estimativas ou evidências centrais. Discuta sinais, magnitudes,
incerteza estatística e relevância substantiva, conforme o método adotado.

\section{Testes de robustez}

Registre especificações alternativas, análises de sensibilidade e outras
verificações que sustentem as conclusões.

\section{Discussão}

Relacione os resultados à pergunta de pesquisa e à literatura apresentada no
capítulo anterior. Explique possíveis divergências e limitações interpretativas.

